\documentclass[10pt,a4paper]{amsart}
\usepackage[margin=19mm]{geometry}
\usepackage{amsmath,amssymb,mathtools}
\usepackage[scr=rsfs,cal=euler]{mathalfa}
\usepackage{xfrac}
\usepackage[unicode,hidelinks,bookmarks=false]{hyperref}
\newcommand{\ie}{i.e.\ }
\newcommand{\cf}{cf.\ }
\newcommand{\resp}{respectively\ }
\newcommand{\Subsection}{\subsection}
\providecommand{\nobreakdash}{\nobreak\hbox{-}}
\setlength{\emergencystretch}{2em}
\multlinegap0pt
\usepackage{enumerate}
\usepackage{amssymb}
\usepackage{tikz}
\usepackage{float}
\newtheorem{lemma}{Lemma}
\title{Shifted L\'evy's Dragon Curve and Directed Graph}
\author{Jonathan Leung}
\date{Source: arXiv:2605.02988v1 (CC BY 4.0)}
\begin{document}
\maketitle

\noindent\emph{Source and licence.} Shifted L\'evy's Dragon Curve and Directed Graph. Version arXiv:2605.02988v1 (CC BY 4.0), distributed by arXiv under the Creative Commons Attribution 4.0 International licence, http://creativecommons.org/licenses/by/4.0/, as recorded in the arXiv licence metadata for this exact version and shown on the abstract page https://arxiv.org/abs/2605.02988v1. Full licence notice: \texttt{licenses/arxiv-2605.02988-CC-BY-4.0.txt}.\par\smallskip

\noindent\textit{Editorial scope.} Counted unit: the article's lemma th:shiftedlevyexplicit (source0.tex lines 445--453). The article's complete original proof of it is delivered with it (source0.tex lines 455--519, one \texttt{proof} environment of 65 lines). No mathematics is written, repaired or re-derived: the statement and the proof are the article's own lines, in the article's own notation, and no retained line is substituted or re-flowed.  Retained uncounted as the source-local context the proof needs: lines 410--419.  Nothing was omitted from any retained span, so the package carries no omission record.  No retained line contains a citation, so the wrapper carries no bibliography and no bibliography entry.  No retained line contains a figure environment, an image-inclusion command or drawing code, so no image is delivered and the package declares no asset dependency.  Editorial formatting only, in addition: an AMS \texttt{amsart} preamble that restates the article's own declarations and macros used by the retained lines, namely the environments \texttt{lemma} on separate counters, together with the standard packages those lines need.  The retained environments print in this document as Lemma 1 in order of appearance, whereas the article prints the same items with the numbering of its own full document.  The complete native source of the article as distributed in the arXiv e-print for this exact version is bundled unchanged as \texttt{sources/199758/source0.tex}.
\begin{equation}
\label{eq:s-functional}
 G(x) =
\begin{cases}
 \alpha G(2x) - \frac{1}{2}
                     &\qquad 0 \leq x < \frac{1}{2}, \\
(1-\alpha) G(2x-1) + \frac{1}{2} & \qquad \frac{1}{2} \leq x \leq 1,\\
 
 \end{cases}
 \end{equation}

\begin{lemma}
\label{th:shiftedlevyexplicit}
For any dyadic point $x \in (0,1)$; that is, there exists $k\in \mathbb{N}$ such that $\omega_k=1$ and $\omega_n(x)=0$ for $\forall n> k$, $G(x)$ has the following expression. 
\begin{equation}
    \label{explicts-levy}
    G(x) = \frac{1}{2}\sum_{n=1}^{k-1}(-1)^{1-\omega_n(x)}\alpha^{n-1-q(x,n-1)}(1-\alpha)^{q(x,n-1)}
\end{equation}
where $\alpha$ = $(1-i)/2$.
\end{lemma} 

\begin{proof}[Proof of Lemma \ref{th:shiftedlevyexplicit}]

Let 
\begin{equation*}
F(x) := \frac{1}{2}\sum_{n=1}^{k-1}(-1)^{1-\omega_n(x)}\alpha^{n-1-q(x,n-1)}(1-\alpha)^{q(x,n-1)}.
\end{equation*}

Recall that $G(x)$ is the unique continuous solution of \eqref{eq:s-functional}.

Let $P(k)$ be the statement: 
$$F(x)=G(x) \mbox{ for all } x \mbox{ of the form } x=(0.\omega_1 \cdots .\omega_{k-1} 1)_{2}.$$

First, consider the base case $k=1$. Since $x=(0.1)_{2}=1/2$, it is clear that the statement $P(1)$ holds as  
\begin{align*}
F(1/2)=F((0.1)_{2}) &= \frac{1}{2}\sum_{n=1}^{0}(-1)^{1- \omega_n(x)}\alpha^{n-1-q(x,n-1)}(1-\alpha)^{q(x,n-1)}\\
& = 0 =G(1/2). 
\end{align*}

Second, consider the case $k=2$, that is $x=(0.01)_2=1/4$ or $x=(0.11)_2= 3/4$

\begin{align*}
F(1/4)=F((0.01)_{2}) &= \frac{1}{2}(-1)^{1-0}\cdot \alpha^0(1-\alpha)^0 =-\frac{1}{2} =G(1/4). \\
F(3/4)=F((0.11)_{2}) &= \frac{1}{2}(-1)^{1-1}\cdot \alpha^0(1-\alpha)^0 =\frac{1}{2} =G(3/4). 
\end{align*}
Therefore the statement P(2) holds.

We continue the proof by induction. Let $l \geq 2$ and assume that the statement $P(l)$ is true; that is,  
$$F(x) = G(x) \mbox { for all } x \mbox{ of the form } x=(0.\omega_1\omega_2\omega_3 \dots \omega_{l-1}1)_2.$$ 
We want to prove $P(l+1)$.

\medskip

Consider $x$ having the form $x=(0.\omega_1\omega_2\omega_3 \dots \omega_{l}1)_2$.
Notice that for $0\leq x<1/2$, $\omega_1=0$ and $q(2x,n-1) = q(x,n)$ and for $1/2 \leq x \leq 1$, $\omega_1=1$ and $q(2x-1,n-1) = q(x,n)-1$.\\
If $0\leq x< 1/2$,
\begin{align*}
    F(x) &= F(0.0 \omega_2\dots\omega_l1)_2\\
    &= \frac{1}{2}\sum_{n=1}^l(-1)^{1-\omega_n(x)}\alpha^{n-1-q(x,n-1)}(1-\alpha)^{q(x,n-1)}\\
    &=-\frac{1}{2}+\frac{1}{2}\sum_{n=2}^l(-1)^{1-\omega_n(x)}\alpha^{n-1-q(x,n-1)}(1-\alpha)^{q(x,n-1)}\\
    &=-\frac{1}{2}+\frac{1}{2}\sum_{n=1}^{l-1}(-1)^{1-\omega_{n+1}(x)}\alpha^{n-q(x,n)}(1-\alpha)^{q(x,n)}\\
    &=-\frac{1}{2}+\frac{\alpha}{2}\sum_{n=1}^{l-1}(-1)^{1-\omega_{n}(2x)}\alpha^{n-1-q(2x,n-1)}(1-\alpha)^{q(2x,n-1)}\\
    &= -\frac{1}{2}+\alpha F(2x)
\end{align*}
Then observe that $2x = (0.\omega2\omega3\omega4\dots\omega_{l-1}1)$.
Therefore $P(l)$ holds by assumption,
\begin{align*}
   F(x) &= -\frac{1}{2}+\alpha G(2x)= G(x)
\end{align*}
If $1/2 \leq x\leq1$,
\begin{align*}
    F(x) &= F(0.1 \omega_2\dots\omega_l1)_2\\
    &= \frac{1}{2}\sum_{n=1}^l(-1)^{1-\omega_n(x)}\alpha^{n-1-q(x,n-1)}(1-\alpha)^{q(x,n-1)}\\
    &=\frac{1}{2}+\frac{1}{2}\sum_{n=2}^l(-1)^{1-\omega_n(x)}\alpha^{n-1-q(x,n-1)}(1-\alpha)^{q(x,n-1)}\\
    &=\frac{1}{2}+\frac{1}{2}\sum_{n=1}^{l-1}(-1)^{1-\omega_{n+1}(x)}\alpha^{n-q(x,n)}(1-\alpha)^{q(x,n)}\\
    &=\frac{1}{2} +\frac{1}{2}\sum_{n=1}^{l-1}(-1)^{1-{\omega_n(2x-1)}}\alpha^{n-q(2x-1,n-1)-1}(1-\alpha)^{q(2x-1,n-1)+1}\\
    &=\frac{1}{2}+\frac{(1-\alpha)}{2}\sum_{n=1}^{l-1}(-1)^{1-\omega_n(2x-1)}\alpha^{n-1-q(2x-1,n-1)}(1-\alpha)^{q(2x-1,n-1)}\\
    &=\frac{1}{2}+(1-\alpha)F(2x-1)
\end{align*}
Again, observe that $2x = (0.\omega2\omega3\omega4\dots\omega_{l-1}1)$.
Therefore $P(l)$ holds by assumption, we have 
\begin{align*}
   F(x) &= \frac{1}{2}+(1-\alpha) G(2x-1)= G(x)
\end{align*}
Therefore, we conclude that the statement $P(k)$ is true for all $k \in \mathbb{N}$.
\end {proof}

\end{document}
